\section{Introduction}\label{sec:intro}

The Agent2Agent (A2A) protocol~\citep{a2a2026} defines how one agent discovers another and hands
it a task. It was contributed to the Linux Foundation, and on August 17, 2026 it joined the
Agentic AI Foundation as a hosted project, which reported backing from more than 150
organizations~\citep{aaifa2a2026}. Teams that need agents from different vendors to exchange work
increasingly have a common transport to build on.

A2A's trust anchor is the Signed Agent Card, which lets a client check that the owner of a
domain issued the card. It does not establish that a peer runs attested, unmodified code; that a
delegating agent actually holds the authority it passes on; that a task payload is confidential
to the intended peer; or that there is an auditable record of who delegated what to whom and
under which policy.

This gap matters once one agent delegates to another. Agent A hands a task to B, and B hands
part of it to C. Nothing in the transport constrains B or C to the authority A actually held. A
card-verified peer can request more than it was granted, and the request will look valid at every
hop. This is the confused deputy problem~\citep{confuseddeputy1988}, now across trust domains
that never negotiated a shared authority model. Recent work treats verifiable delegation for
agents as open: Prakash reports finding no prior implemented protocol that jointly combined
verifiable delegation, holder-side attenuation, chained policy, transport bindings, and provenance
records, and proposes one~\citep{aip2026}; Forough et al.\ list compound attestation for multi-hop
agent chains among the open challenges for confidential agentic AI~\citep{secretssurvey2026}.
Mutual TLS and OAuth scopes secure the connection and assert an identity; they do not attenuate
authority per hop, appraise runtime integrity, or seal a payload to an appraised peer.

This report describes cA2A (Confidential A2A), a trust \emph{profile} layered on A2A rather than a
competing transport. cA2A does not move tasks between agents; A2A does that. It defines the trust
envelope around a delegated task, in the way TRACE~\citep{trace2025} binds to
RATS~\citep{rats2021} and EAT~\citep{eat2023} rather than replacing them. cA2A composes four
mechanisms:

\begin{itemize}[leftmargin=1.5em]
  \item \textbf{Attenuated delegation.} Each hop carries a signed delegation credential whose
    scope is a subset of its parent's. Depth is bounded, the root issuer must be trusted by the
    verifier, and a credential identifier cannot repeat within a chain.
  \item \textbf{Runtime attestation.} A caller appraises a peer's attestation evidence, which
    binds the peer's channel key to a measurement, before it sends a delegated task.
  \item \textbf{Sealed peer channel.} The task payload is sealed to the channel key the appraised
    evidence vouches for, so only the holder of the matching private key can open it.
  \item \textbf{Provenance record.} Each hop emits a signed TRACE record referencing its parent
    record's hash and the delegation credential it acted under, so an auditor can check the
    path offline.
\end{itemize}

None of these mechanisms is new in isolation. Macaroons introduced offline attenuation of bearer
credentials~\citep{macaroons2014}, and recent agent-specific work specifies per-hop attenuation
with Cedar policies and parent-swap resistance~\citep{pedigree2026} or implements attenuated,
provenance-bound capability tokens across MCP and A2A~\citep{aip2026}. Section~\ref{sec:related}
compares them.

\subsection*{Contributions}

\begin{itemize}[leftmargin=1.5em]
  \item \textbf{The cA2A profile.} A binding on A2A, and on the IETF RATS, EAT, and
    SCITT~\citep{scitt2024} work, that composes attenuated delegation, attestation-gated sealing,
    and signed provenance for agent-to-agent calls without changing the A2A transport.
  \item \textbf{An open reference implementation.} The delegation credential model and offline
    chain verifier, the inbound decision path, the sealed channel, attestation verifiers for AMD
    SEV-SNP, Intel TDX, and TPM~2.0 evidence, signed TRACE provenance, and a CLI, released under
    the MIT license as ca2a-runtime~\citep{ca2arepo2026}. Delegation semantics mirror those of
    agent-manifest~\citep{agentmanifest2026}, and hardware-evidence parsing and certificate-chain
    checks are shared with it.
  \item \textbf{Six properties with their evidence stated.} We state six properties
    (Section~\ref{sec:claims}) and report software experiments rerun against ca2a 0.3.1
    (Section~\ref{sec:evaluation}). C1 (attenuation soundness) and C2 (in-chain replay and splice
    rejection) are validated in software. C3 (scope-policy intersection) is validated in software,
    including the inbound path of the reference transport. C4 (sealed-payload confidentiality) is
    validated at the cryptographic layer only. C5 (provenance integrity) is validated for
    structural consistency by the experiment, while authenticity rests on signed records covered by
    the project's tests. C6 (cross-operator attestation) is exercised with synthetic evidence and
    is not validated on hardware across independent operators.
\end{itemize}

\subsection*{Paper structure}

Section~\ref{sec:background} covers A2A and its Signed Agent Card, TEE attestation, the TRACE
record format, Cedar policy, and provenance. Section~\ref{sec:design} describes the cA2A profile
and how its mechanisms compose. Section~\ref{sec:claims} states the six properties and the status
of each. Section~\ref{sec:evaluation} presents the rerun experiments. Section~\ref{sec:after}
reports findings made after the original evaluation. Section~\ref{sec:related} surveys related
work, and Section~\ref{sec:conclusion} concludes.
